\SourceLeaf{176}Es ist leicht zu beweisen, dass, wenn Euklids Geometrie nicht die wahre ist,
es gar keine ähnlichen Figuren gibt. Die Winkel in einem gleichseitigen Dreieck
sind dann auch je nach der Länge der Seite verschieden, wobei ich gar nichts
Absurdes finde. Der Winkel ist dann eine Funktion der Seite und die Seite eine Funktion
des Winkels, natürlich eine solche Funktion, in der zugleich eine
konstante Linie vorkommt. Es scheint etwas paradox, dass eine konstante Linie
  gleichsam a priori möglich sein könne. Ich finde aber darin nichts
Widersprechendes. Es wäre sogar wünschenswert, dass die Geometrie Euklids nicht
  wahr wäre, weil wir dann ein allgemeines Maß a priori hätten, z. B.
könnte man als Raumeinheit die Seite desjenigen gleichseitigen Dreiecks annehmen,
  dessen Winkel $=59^{0} 59' 59''{,}99999$. .\ .\ .\ .\ .\ .\ .

\begin{center}\rule{3cm}{0.4pt}\end{center}
